% Presentation of original arrow notation. See flow-manifest.json.
\begin{center}
\begin{tikzpicture}[book flow node/.style={draw=BookHeading,line width=.65pt,fill=BookPanel,rounded corners=1.5pt,align=center,inner sep=5pt,minimum height=10mm,font=\fontsize{9}{11}\selectfont,execute at begin node={\hyphenpenalty=10000\exhyphenpenalty=10000}},book flow edge/.style={-{Latex[length=2mm]},draw=BookHeading,line width=.7pt}]
\node[book flow node,text width=48mm] (n0) at (-3,0) {BIJOU\_\allowbreak{}ACCESSIBLE=1};
\node[book flow node,text width=48mm] (n1) at (3,0) {'accessible'};
\node[book flow node,text width=48mm] (n2) at (-3,-2) {TERM=dumb};
\node[book flow node,text width=48mm] (n3) at (3,-2) {'pipe'};
\node[book flow node,text width=48mm] (n4) at (-3,-4) {!stdout.\allowbreak{}is\allowbreak{}TTY};
\node[book flow node,text width=48mm] (n5) at (3,-4) {'pipe'};
\node[book flow node,text width=48mm] (n6) at (-3,-6) {CI};
\node[book flow node,text width=48mm] (n7) at (3,-6) {'static'};
\node[book flow node,text width=48mm] (n8) at (-3,-8) {stdout.\allowbreak{}is\allowbreak{}TTY};
\node[book flow node,text width=48mm] (n9) at (3,-8) {'interactive'};
\draw[book flow edge] (n0) -- (n1);
\draw[book flow edge] (n2) -- (n3);
\draw[book flow edge] (n4) -- (n5);
\draw[book flow edge] (n6) -- (n7);
\draw[book flow edge] (n8) -- (n9);
\end{tikzpicture}
\end{center}
